\section{Related Work}
\label{sec:related}

We discuss related work on eBPF program optimization, operation extensions, compiler correctness, and kernel extension safety. We compare how optimizations and extensions are integrated into eBPF and how program safety and semantic preservation are established.

\noindent\textbf{eBPF program optimization.} Existing eBPF optimizations differ in where programs are transformed and how optimization results are represented. K2~\cite{xu2021synthesizing}, Merlin~\cite{mao2024merlin}, and EPSO~\cite{zhu2025epso} rewrite programs in userspace, producing bytecode in the existing eBPF instruction set. KFuse~\cite{kfuse} merges separately verified programs in the kernel, while ePass~\cite{xiang2025epass} provides a program transformation framework that cooperates with the verifier. Shahinfar et al.~\cite{shahinfar2026fun} propose post-verification optimization (PVO), in which a trusted userspace optimizer compiles verified eBPF programs into native code for kernel execution. Compared with bytecode rewriting within the existing instruction set, \tool uses \kinsns to select additional native implementations. Unlike PVO's userspace code generation, \tool checks \kinsn use and generates native code in the kernel, without trusting the userspace pattern recognizer for kernel safety.

\noindent\textbf{Extending eBPF operations.} eBPF extensions differ in how they represent and check new operations. kfuncs~\cite{kfuncs} expose kernel functions through a call interface. The verifier checks usage conditions, while function bodies remain outside bytecode analysis. hXDP~\cite{brunella2020hxdp} extends the eBPF instruction set and execution engine for FPGAs and uses hardware checks to constrain execution. Both \tool and hXDP provide additional functionality through extended instructions, but \tool targets kernel execution on CPUs and submits new operations' eBPF expansions to the existing verifier.

\noindent\textbf{Compiler correctness.} Research on verifying and synthesizing compilers for kernel domain-specific languages focuses on semantic preservation during translation. Jitk~\cite{wang2014jitk} verifies compilation paths, Jitterbug~\cite{nelson2020specification} checks BPF JIT translation steps, and JitSynth~\cite{vangeffen2020jitsynth} synthesizes correct JIT compilation fragments. Our correctness argument for \kinsns likewise requires semantic preservation: the proof concerns the correspondence between their eBPF expansions and native implementations (Section~\ref{sec:formal-verification}).

\noindent\textbf{Kernel extension safety.} Research on kernel extension safety uses mechanisms such as external proofs, language restrictions, and runtime protection. Proof-carrying code (PCC)~\cite{necula1997pcc} requires consumers to check safety proofs accompanying code. BCF~\cite{sun2025prove} improves eBPF verifier precision using external proofs checked in the kernel. KFlex~\cite{dwivedi2024fast} combines static checking of kernel interfaces with runtime protection. Rex~\cite{jia2025rex} replaces a separate bytecode verifier with a restricted subset of safe Rust, a trusted toolchain, and a runtime. MOAT~\cite{lu2024moat} and SafeBPF~\cite{soo2024safebpf} retain the verifier and add runtime memory isolation.